%==============================================================================
%  thesis.tex  --  Documento raíz de la tesis doctoral
%  Universidad Nacional de San Luis (UNSL)
%
%  Cómo compilar:
%    Local:    latexmk thesis.tex        (configuración en .latexmkrc)
%    Overleaf: compilador pdfLaTeX; minted ya funciona (shell-escape on).
%
%  Estructura del proyecto:
%    config/     paquetes (packages), macros (commands) y datos (metadata)
%    styles/     estilos heredados — rac.sty es el FORMATO EXIGIDO por la UNSL
%    chapters/   capítulos 00_resumen ... 08_conclusiones
%    appendices/ apéndices A, B, ...
%    images/     figuras por capítulo        tables/  tablas por capítulo
%    build/      salida de la compilación (no se versiona)
%==============================================================================
\documentclass[12pt,oneside,reqno]{report}

% Buscar paquetes/estilos propios también en styles/
\makeatletter
\def\input@path{{styles/}}
\makeatother

% rac.sty define el formato de tesis (portada, márgenes, índices). Se carga
% temprano —antes que hyperref— para preservar los marcadores del PDF.
\usepackage{rac}

%==============================================================================
%  config/packages.tex  --  Paquetes y configuracion base
%  Tesis doctoral - Universidad Nacional de San Luis (UNSL)
%------------------------------------------------------------------------------
%  Preambulo modernizado y deduplicado (junio 2026).
%  Cambios clave respecto del template original:
%    - Se ELIMINO el conflicto color/xcolor: ahora se carga SOLO xcolor.
%    - Se quitaron cargas DUPLICADAS (graphicx, amsmath, setspace, color...).
%    - Se removieron paquetes innecesarios para esta tesis (epsfig, CJKutf8,
%      fontenc LGR/griego).
%    - Se agregaron microtype y booktabs (tipografia y tablas modernas).
%    - hyperref se carga casi al final, como corresponde, y cleveref despues.
%    - Los paquetes de algoritmos se cargan de forma condicional.
%==============================================================================

% ---- Codificacion e idioma --------------------------------------------------
\usepackage[utf8]{inputenc}      % archivos fuente en UTF-8 (acentos, n)
\usepackage[T1]{fontenc}         % salida con acentos correctos / copiar-pegar
\usepackage[spanish,es-tabla,es-noindentfirst,es-nodecimaldot]{babel}
% es-tabla: los flotantes dicen "Tabla"; es-noindentfirst: sin sangria inicial.

% ---- Fuentes ----------------------------------------------------------------
% Tipografia: Computer Modern (la que usa la tesis de referencia de la UNSL).
% A proposito NO se carga lmodern, para que la fuente coincida con la del
% director (con T1 + cm-super salen las mismas SFRM*; en TeX Live completo y
% Overleaf ya estan). Si alguna vez quisieras la version moderna, agrega:
%   \usepackage{lmodern}
\usepackage{courier}             % monoespaciada para codigo en linea
% microtype: protrusion (mejora margenes). expansion=false para no depender de
% fuentes escalables (cm-super) al usar Computer Modern -> compila en todos lados.
\usepackage[expansion=false]{microtype}

% ---- Matematica -------------------------------------------------------------
\usepackage{amsmath}
\usepackage{amssymb}
\usepackage{amsfonts}
\usepackage{amsthm}
\usepackage{amsxtra}
\usepackage{mathtools}           % extiende amsmath (\DeclarePairedDelimiter...)

% ---- Color: UN SOLO paquete (xcolor reemplaza por completo a color) ----------
\usepackage[dvipsnames,svgnames,table]{xcolor}
% 'table' habilita colores de fila/celda; 'dvipsnames'/'svgnames' dan nombres.
% NO cargar 'color' por separado: ese era el origen de los errores de choque.

% Colores institucionales UNSL (del escudo), para acentos del template:
\definecolor{unslazul}{RGB}{48,48,144}       % azul del escudo
\definecolor{unslceleste}{RGB}{120,192,240}  % celeste del escudo

% ---- Figuras y tablas -------------------------------------------------------
\usepackage{graphicx}
\graphicspath{{images/}}         % raiz para \includegraphics (mas subcarpetas)
\usepackage{rotating}            % figuras/tablas apaisadas
% rac.sty redefine \@makecaption con un formato propio que el paquete caption
% no reconoce (warning "Unknown document class") y que igualmente pisa con sus
% defaults. Restauramos la definicion estandar de LaTeX antes de cargarlo para
% que la reconozca; el resultado visual es el mismo (manda caption).
\makeatletter
\long\def\@makecaption#1#2{%
  \vskip\abovecaptionskip
  \sbox\@tempboxa{#1: #2}%
  \ifdim \wd\@tempboxa >\hsize
    #1: #2\par
  \else
    \global \@minipagefalse
    \hb@xt@\hsize{\hfil\box\@tempboxa\hfil}%
  \fi
  \vskip\belowcaptionskip}
\makeatother
\usepackage{subcaption}          % subfiguras (reemplaza al viejo 'subfigure')
\usepackage[font={small,bf}]{caption}
\usepackage{booktabs}            % \toprule \midrule \bottomrule
\usepackage{multirow}
\usepackage{makecell}
\usepackage{adjustbox}           % escalar/ajustar tablas y cajas anchas

% ---- Espaciado, notas, varios -----------------------------------------------
\usepackage{setspace}            % \onehalfspacing, \doublespacing
\usepackage[bottom]{footmisc}    % notas al pie siempre al fondo de la pagina
\usepackage{verbatim}
\usepackage{soul}                % \hl (resaltar), \st (tachar)
\usepackage{etoolbox}            % \patchcmd y utilidades de programacion

% ---- Pagina: A4, margenes simetricos y encabezado ---------------------------
% Reemplaza los margenes "Rackham" de rac.sty por A4 con margenes simetricos
% (igual izquierda/derecha) y mas aire arriba/abajo. Agrega un encabezado con
% el nombre del capitulo y el numero de pagina, con una linea debajo.
\usepackage[a4paper,left=3cm,right=3cm,top=3cm,bottom=3cm,
            headheight=15pt,headsep=0.7cm]{geometry}
\usepackage{fancyhdr}
\pagestyle{fancy}
\fancyhf{}                                  % limpia encabezado y pie
\fancyhead[L]{\nouppercase{\leftmark}}      % capitulo a la izquierda
\fancyhead[R]{\thepage}                     % numero de pagina a la derecha
\renewcommand{\headrulewidth}{0.4pt}        % linea bajo el encabezado
\renewcommand{\footrulewidth}{0pt}
% Aperturas de capitulo: sin encabezado, numero de pagina abajo-centro
\fancypagestyle{plain}{%
  \fancyhf{}%
  \fancyfoot[C]{\thepage}%
  \renewcommand{\headrulewidth}{0pt}%
}

% ---- Algoritmos (carga condicional) -----------------------------------------
% algorithm + algorithmicx/algpseudocode estan en TeX Live completo y Overleaf.
% Se cargan solo si estan disponibles para no romper entornos minimos.
\IfFileExists{algpseudocode.sty}{%
  \usepackage{algorithm}%
  \usepackage{algpseudocode}%
}{%
  \PackageWarning{packages}{algpseudocode no encontrado: se omiten los entornos de algoritmos.}%
}

% ---- Codigo fuente: minted --------------------------------------------------
% minted requiere compilar con  -shell-escape  y tener Pygments instalado
% (pip install Pygments). En Overleaf funciona sin configuracion extra.
% Alternativa sin dependencias: comentar esta linea y usar 'listings'.
% La opcion outputdir ya no es necesaria (ni valida) con minted v3+ / TeX Live 2024+:
% el directorio de salida se detecta automaticamente.
\usepackage{minted}

% ---- Bibliografia -----------------------------------------------------------
\usepackage{natbib}              % citas \citep / \citet (estilo en thesis.tex)
% Por defecto se usa plainnat. El template trae agu04.bst en styles/ como
% alternativa (estilo de la AGU); para usarlo: \bibliographystyle{agu04}.
% El encabezado de la bibliografia lo pone \startbibliography (estilo capitulo);
% anulamos el encabezado propio de natbib para que no quede duplicado:
\renewcommand{\bibsection}{}

% ---- Abreviaturas / acronimos -----------------------------------------------
\usepackage[printonlyused]{acronym}   % lista de abreviaturas (lee abbr.tex)

% ---- Hiperenlaces (cargar casi al final) ------------------------------------
\usepackage[spanish,unicode]{hyperref}
\hypersetup{
  colorlinks=true,
  linkcolor=unslazul,         % enlaces internos (capitulos, secciones)
  citecolor=unslazul,         % citas bibliograficas
  urlcolor=unslazul,          % URLs
  linktoc=all,                % en el indice, enlaza numero y texto
}
\renewcommand{\UrlFont}{\color{unslazul}\rmfamily}

% ---- Referencias cruzadas inteligentes (despues de hyperref) ----------------
\usepackage[capitalise,noabbrev]{cleveref}   % \cref{...} -> "figura 3.1", etc.

% rac.sty usa un tipo de entrada 'chap' en el indice; le indicamos a hyperref
% su nivel de marcador para evitar el warning "bookmark level for unknown chap".
\makeatletter
\providecommand*{\toclevel@chap}{0}
\makeatother

% ---- Nota sobre estilos heredados del template ------------------------------
% styles/aas_macros.sty  : macros de revistas astronomicas (ADS). NO se carga
%   porque no es necesario en esta tesis; queda por compatibilidad. Si lo
%   necesitaras: \usepackage{aas_macros}
% styles/agu04.bst        : estilo bibliografico de la AGU (geofisica), opcional.
% styles/rac.sty          : formato de tesis exigido (UNSL/Rackham); SE USA.

    % paquetes y configuración (hyperref al final)
%==============================================================================
%  config/commands.tex  --  Colores, entornos y macros propias
%  Tesis doctoral - Universidad Nacional de San Luis (UNSL)
%==============================================================================

% ---- Colores ----------------------------------------------------------------
\definecolor{codebg}{rgb}{0.97,0.97,0.97}   % fondo de los bloques de codigo
\definecolor{codeframe}{rgb}{0.80,0.80,0.80}

% ---- Configuracion de minted (codigo fuente) --------------------------------
% Se define aca porque depende de los colores de arriba. minted ya fue cargado
% en config/packages.tex.
\setminted{
  fontsize=\small,
  breaklines=true,        % corta lineas largas
  breakanywhere=true,
  tabsize=4,
  frame=lines,
  framesep=2mm,
  bgcolor=codebg,
  linenos=false,
}
\newcommand{\code}[1]{\mintinline{text}{#1}}   % codigo corto en linea

% ---- Referencias cruzadas en espanol (cleveref) -----------------------------
\crefname{figure}{figura}{figuras}
\Crefname{figure}{Figura}{Figuras}
\crefname{table}{tabla}{tablas}
\Crefname{table}{Tabla}{Tablas}
\crefname{equation}{ecuacion}{ecuaciones}
\Crefname{equation}{Ecuacion}{Ecuaciones}
\crefname{chapter}{capitulo}{capitulos}
\Crefname{chapter}{Capitulo}{Capitulos}
\crefname{section}{seccion}{secciones}
\Crefname{section}{Seccion}{Secciones}
\crefname{algorithm}{algoritmo}{algoritmos}
\Crefname{algorithm}{Algoritmo}{Algoritmos}

% ---- Operadores y delimitadores matematicos ---------------------------------
\DeclareMathOperator*{\argmax}{arg\,max}
\DeclareMathOperator*{\argmin}{arg\,min}
\DeclarePairedDelimiter{\ceil}{\lceil}{\rceil}
\DeclarePairedDelimiter{\floor}{\lfloor}{\rfloor}
\DeclarePairedDelimiter{\abs}{\lvert}{\rvert}

% ---- Entornos de teoremas (numerados por capitulo, en espanol) --------------
\theoremstyle{plain}
\newtheorem{teorema}{Teorema}[chapter]
\newtheorem{proposicion}[teorema]{Proposicion}
\newtheorem{corolario}[teorema]{Corolario}
\newtheorem{lema}[teorema]{Lema}
\newtheorem{conjetura}[teorema]{Conjetura}

\theoremstyle{definition}
\newtheorem{definicion}[teorema]{Definicion}
\newtheorem{ejemplo}[teorema]{Ejemplo}

\theoremstyle{remark}
\newtheorem{observacion}[teorema]{Observacion}
\newtheorem{nota}[teorema]{Nota}

\renewcommand{\proofname}{Demostraci\'on}   % entorno proof en espanol

% rac.sty pone el numero de capitulo en romano; forzamos numeracion arabe
% para que los teoremas queden "1.2" en lugar de "I.2" (consistente con secciones).
\renewcommand{\theteorema}{\arabic{chapter}.\arabic{teorema}}

% ---- Palabras clave de algoritmos en espanol (solo si el paquete esta) ------
\makeatletter
\@ifpackageloaded{algpseudocode}{%
  \floatname{algorithm}{Algoritmo}%
  \renewcommand{\algorithmicrequire}{\textbf{Entrada:}}%
  \renewcommand{\algorithmicensure}{\textbf{Salida:}}%
  \renewcommand{\algorithmicif}{\textbf{si}}%
  \renewcommand{\algorithmicthen}{\textbf{entonces}}%
  \renewcommand{\algorithmicelse}{\textbf{si no}}%
  \renewcommand{\algorithmicend}{\textbf{fin}}%
  \renewcommand{\algorithmicwhile}{\textbf{mientras}}%
  \renewcommand{\algorithmicdo}{\textbf{hacer}}%
  \renewcommand{\algorithmicfor}{\textbf{para cada}}%
  \renewcommand{\algorithmicforall}{\textbf{para todo}}%
  \renewcommand{\algorithmicreturn}{\textbf{retornar}}%
  \renewcommand{\algorithmicprocedure}{\textbf{procedimiento}}%
  \renewcommand{\algorithmicfunction}{\textbf{funcion}}%
  \algtext*{EndWhile}% quita el "fin mientras"
  \algtext*{EndIf}%    quita el "fin si"
  \algtext*{EndFor}%   quita el "fin para"
}{}
\makeatother

% ---- Marca de pendiente para borradores -------------------------------------
% Uso:  \todo{revisar esta cifra}
\newcommand{\todo}[1]{%
  \textbf{\textcolor{red}{[TODO: #1]}}%
  \marginpar{\textcolor{red}{\tiny$\blacktriangleleft$ TODO}}%
}

% ---- Entorno de dedicatoria -------------------------------------------------
\newenvironment{dedicatoria}
  {\clearpage \thispagestyle{empty} \vspace*{\stretch{1}} \slshape \raggedleft}
  {\par \vspace{\stretch{3}} \clearpage}

% ---- Encabezado para secciones preliminares sin número (Resumen, Abstract) --
% Replica el estilo centrado de rac.sty y lo agrega al índice.
\newcommand{\frontsection}[1]{%
  \newpage
  \addcontentsline{toc}{chapter}{#1}%
  \hbox{}\twoinmar
  \begin{center}\large\bfseries #1\end{center}%
  \vspace{0.4in}%
}
    % colores, minted, teoremas, algoritmos, macros
%==============================================================================
%  config/metadata.tex  --  Datos de la tesis y portada
%  >>> EDITA ESTOS VALORES <<<
%==============================================================================

% ---- Datos principales ------------------------------------------------------
\newcommand{\thesisTitle}{Integración de Modelos de Lenguaje y planificación a partir de flujos conversacionales inducidos en agentes híbridos para sistemas de diálogo orientado a tareas}
\newcommand{\thesisAuthor}{Mg. Juan Manuel Fernández}
\newcommand{\thesisProgram}{Doctorado en Ciencias de la Computación}
\newcommand{\thesisFaculty}{Facultad de Ciencias Físico, Matemáticas y Naturales}
\newcommand{\thesisUniversity}{Universidad Nacional de San Luis}
\newcommand{\thesisYear}{2026}
\newcommand{\thesisLogo}{logo_unsl}        % archivo en images/

% ---- Dirección --------------------------------------------------------------
\newcommand{\thesisDirector}{Dr. Sergio Burdisso, Director}
\newcommand{\thesisCoDirector}{Dr. Marcelo Errecalde, Co-Director}

% ---- Metadatos del PDF (hyperref ya está cargado) ---------------------------
\hypersetup{
  pdftitle={\thesisTitle},
  pdfauthor={\thesisAuthor},
  pdfsubject={Tesis de \thesisProgram},
}

% ---- Portada -----------------------------------------------------------------
% Diseño propio (inspirado en la tesis de Maestría): título en color con línea,
% bloque Autor/Dirección a dos columnas, logo y datos, con buen espaciado.
% Para una portada más sobria, cambiá 'unslazul' por 'black'.
\colorlet{portada}{unslazul}            % azul del escudo UNSL (definido en packages.tex)

\newcommand{\maketitlepage}{%
  \thispagestyle{empty}%
  \begingroup
  \centering
  \vspace*{\stretch{1.3}}
  {\Large\bfseries\color{portada}\thesisTitle\par}   % titulo (mas chico que antes)
  \vspace{1em}
  {\color{portada}\rule{\linewidth}{1pt}}\par
  \vspace{\stretch{1}}
  % --- Autor / Dirección, a dos columnas ---
  \noindent
  \begin{minipage}[t]{0.46\linewidth}
    \raggedright\itshape Autor:\par\smallskip
    \upshape\large \thesisAuthor
  \end{minipage}\hfill
  \begin{minipage}[t]{0.46\linewidth}
    \raggedleft\itshape Dirección:\par\smallskip
    \upshape \thesisDirector\\
    \thesisCoDirector
  \end{minipage}\par
  \vspace{16mm}                                       % separacion fija: directores no se pegan al grado
  Tesis presentada para obtener el grado de\par\smallskip
  {\large\bfseries \thesisProgram}\par
  \vspace{\stretch{1.3}}
  \includegraphics[width=3.2cm]{\thesisLogo}\par
  \vspace{\stretch{1.3}}
  \thesisFaculty\par\smallskip
  {\large\color{portada}\thesisUniversity}\par\smallskip
  San Luis, Argentina\par
  \vspace{1.4em}
  \thesisYear\par
  \vspace*{\stretch{1.3}}
  \endgroup
  \clearpage
}
    % título, autor, director, año...

\begin{document}

% Estilo bibliográfico (natbib). Alternativas: abbrvnat, unsrtnat, agu04.
\bibliographystyle{plainnat}

%--------------------------------- PORTADA ------------------------------------
\maketitlepage
\newpage\null\thispagestyle{empty}\newpage

%----------------------------- MATERIAL PRELIMINAR ----------------------------
\initializefrontsections     % numeración romana (i, ii, iii...)

% Opcionales (descomentar para usarlos):
% \begin{dedicatoria} A quien corresponda. \end{dedicatoria}
% \startacknowledgementspage %========================== Agradecimientos (opcional) ========================
% Se incluye descomentando en thesis.tex:
%   \startacknowledgementspage %========================== Agradecimientos (opcional) ========================
% Se incluye descomentando en thesis.tex:
%   \startacknowledgementspage %========================== Agradecimientos (opcional) ========================
% Se incluye descomentando en thesis.tex:
%   \startacknowledgementspage \input{chapters/agradecimientos}
\todo{Escribir los agradecimientos.}

\todo{Escribir los agradecimientos.}

\todo{Escribir los agradecimientos.}


%======================= Resumen (material preliminar) ========================
\frontsection{Resumen}

\todo{Escribir el resumen (200--400 palabras): problema abordado, enfoque
propuesto, principales contribuciones y resultados.}

\bigskip
\noindent\textbf{Palabras clave:} palabra1, palabra2, palabra3, palabra4.

% --- Versión en inglés (opcional, habitual en doctorados argentinos) ---------
\frontsection{Abstract}

\todo{Write the English abstract here (a translation of the Spanish summary).}

\bigskip
\noindent\textbf{Keywords:} keyword1, keyword2, keyword3, keyword4.
  % Resumen (y Abstract)

\tableofcontents
\listoffigures             % descomentar si hay más de una figura
\listoftables              % descomentar si hay más de una tabla
\listofabbreviations         % lista de abreviaturas (lee abbr.tex)
\newpage\null\thispagestyle{empty}\newpage

%--------------------------------- CAPÍTULOS ----------------------------------
\startthechapters            % numeración arábiga y capítulos numerados

\chapter{Introducción}\label{cap:introduccion}
%=============================== Introducción =================================
% Estructura sugerida: contexto y motivación, planteo del problema, objetivos,
% contribuciones y organización del documento.

La \ac{IA} y, en particular, el \ac{PLN} han transformado la manera en que
se procesa la información textual. \todo{Contextualizar el área de la tesis.}

\section{Motivación y planteo del problema}
\todo{Describir el problema de investigación y por qué es relevante. Plantear
la brecha (gap) que la tesis viene a cubrir.}

\section{Objetivos}
\subsection{Objetivo general}
\todo{Enunciar el objetivo general en una o dos oraciones.}

\subsection{Objetivos específicos}
\begin{itemize}
  \item \todo{Objetivo específico 1.}
  \item \todo{Objetivo específico 2.}
  \item \todo{Objetivo específico 3.}
\end{itemize}

\section{Contribuciones}
Las principales contribuciones de esta tesis son:
\begin{enumerate}
  \item \todo{Contribución 1.}
  \item \todo{Contribución 2.}
\end{enumerate}

\section{Organización de la tesis}
El resto del documento se organiza como sigue. El \cref{cap:estado-del-arte}
revisa los antecedentes; el \cref{cap:marco-teorico} presenta el marco teórico;
el \cref{cap:metodologia} describe la metodología; los
\cref{cap:experimentos,cap:resultados} detallan los experimentos y sus
resultados; el \cref{cap:discusion} los discute; y el \cref{cap:conclusiones}
cierra con las conclusiones y el trabajo futuro.


\chapter{Estado del arte}\label{cap:estado-del-arte}
%=============================== Estado del arte ==============================
\todo{Revisión crítica de la literatura relevante. Agrupar por enfoques o
líneas de trabajo, no por autor.}

\section{Trabajos relacionados}
A modo de ejemplo, \citet{ejemplo2023} propusieron un enfoque basado en X,
mientras que otros trabajos abordan el problema desde Y \citep{ejemplo2024}.
\todo{Reemplazar por las referencias reales (cargarlas en \texttt{references.bib}).}

\section{Síntesis y brechas}
\todo{Resumir el estado del arte e identificar explícitamente las brechas que
esta tesis aborda.}


\chapter{Marco teórico}\label{cap:marco-teorico}
%=============================== Marco teórico ================================
\todo{Fundamentos teóricos necesarios para comprender la tesis.}

\section{Conceptos preliminares}
\begin{definicion}\label{def:ejemplo}
Una \emph{definición} se escribe con el entorno \texttt{definicion} y queda
numerada por capítulo (esta es la \cref{def:ejemplo}).
\end{definicion}

Las ecuaciones numeradas se referencian con la \cref{eq:modelo}:
\begin{equation}\label{eq:modelo}
  f(\mathbf{x}) = \sum_{i=1}^{n} w_i\, x_i + b .
\end{equation}

\begin{teorema}\label{teo:ejemplo}
Enunciado de un teorema de ejemplo.
\end{teorema}

\begin{proof}
\todo{Demostración (o quitar el entorno si no aplica).}
\end{proof}

\section{...}
\todo{Desarrollar el resto del marco teórico.}


\chapter{Metodología}\label{cap:metodologia}
%=============================== Metodología ==================================
\todo{Describir el enfoque metodológico y los materiales utilizados.}

\section{Enfoque propuesto}
\todo{Explicar la propuesta paso a paso.}

% --- Ejemplo de figura -------------------------------------------------------
% Poné la imagen en images/04_metodologia/ y descomentá:
%
% \begin{figure}[htbp]
%   \centering
%   \includegraphics[width=0.7\textwidth]{04_metodologia/diagrama}
%   \caption{Esquema general del enfoque propuesto.}
%   \label{fig:enfoque}
% \end{figure}

% --- Ejemplo de algoritmo ----------------------------------------------------
% Requiere los paquetes algorithm/algpseudocode (presentes en TeX Live completo
% y en Overleaf). Descomentá para usarlo:
%
% \begin{algorithm}[htbp]
%   \caption{Procedimiento de ejemplo}\label{alg:ejemplo}
%   \begin{algorithmic}[1]
%     \Require conjunto de entrada $X$
%     \Ensure resultado $Y$
%     \For{cada $x \in X$}
%       \If{$x$ cumple la condición}
%         \State procesar $x$
%       \EndIf
%     \EndFor
%     \State \Return $Y$
%   \end{algorithmic}
% \end{algorithm}

\section{Materiales}
\todo{Datos, herramientas y recursos empleados.}


\chapter{Experimentos}\label{cap:experimentos}
%=============================== Experimentos =================================
\todo{Diseño experimental: datos, configuración, métricas y baselines.}

\section{Conjuntos de datos}
\todo{Describir los datos utilizados.}

\section{Configuración experimental}
Un fragmento de código de ejemplo, resaltado con \texttt{minted}:

\begin{minted}{python}
def accuracy(y_true, y_pred):
    """Proporción de predicciones correctas."""
    correctas = sum(int(a == b) for a, b in zip(y_true, y_pred))
    return correctas / len(y_true)
\end{minted}

\section{Métricas}
Las métricas reportadas se resumen en la \cref{tab:metricas}.

\begin{table}[htbp]
  \centering
  \caption{Resultados de ejemplo sobre el conjunto de prueba.}
  \label{tab:metricas}
  \begin{tabular}{lcc}
    \toprule
    Modelo            & Precisión & F\textsubscript{1} \\
    \midrule
    Línea base        & 0,71      & 0,69 \\
    Enfoque propuesto & \textbf{0,84} & \textbf{0,83} \\
    \bottomrule
  \end{tabular}
\end{table}


\chapter{Resultados}\label{cap:resultados}
%================================ Resultados ==================================
\todo{Presentar los resultados de forma objetiva (la interpretación va en la
Discusión).}

\section{Resultados principales}
\todo{Tablas y figuras con los resultados centrales.}

\section{Análisis de sensibilidad}
\todo{Variaciones, ablations, robustez.}


\chapter{Discusión}\label{cap:discusion}
%================================ Discusión ===================================
\todo{Interpretar los resultados a la luz de los objetivos y del estado del arte.}

\section{Interpretación de los resultados}
\todo{¿Qué significan los resultados? ¿Confirman o refutan la hipótesis?}

\section{Limitaciones}
\todo{Limitaciones del estudio y amenazas a la validez.}


\chapter{Conclusiones y trabajo futuro}\label{cap:conclusiones}
%======================== Conclusiones y trabajo futuro =======================
\section{Conclusiones}
\todo{Sintetizar las contribuciones y los hallazgos principales, en relación
con los objetivos planteados en la Introducción.}

\section{Trabajo futuro}
\todo{Líneas de investigación que quedan abiertas.}


%--------------------------------- APÉNDICES ----------------------------------
\startappendices
\appendix{Material suplementario}\label{ap:material-suplementario}
%========================= Apéndice A — Material suplementario ================
\todo{Material complementario: derivaciones extensas, pruebas, detalles de
implementación, hiperparámetros, etc.}


\appendix{Resultados adicionales}\label{ap:resultados-adicionales}
%========================= Apéndice B — Resultados adicionales ================
\todo{Tablas y figuras complementarias que no entran en el cuerpo principal.}


%------------------------------- BIBLIOGRAFÍA ---------------------------------
\startbibliography
\newpage\null\thispagestyle{empty}\newpage
\begin{singlespace}
  \bibliography{references}
\end{singlespace}

\end{document}
